\chapter{Construcción y Enriquecimiento del Corpus MetaSMS-HSS}
\label{dataset}

La calidad, trazabilidad y coherencia de los datos constituyen el pilar fundamental de cualquier sistema de clasificación de texto. En este capítulo se detalla la construcción de \textbf{MetaSMS-HSS} (\textit{Meta-dataset Unificado de SMS para clasificación Ham/Spam/Smishing}), un corpus maestro diseñado no solo para maximizar el volumen de muestras, sino para constituir una herramienta unificada, trazable y rigurosa de cara a la experimentación. Se aborda la detección de tres clases, poniendo especial énfasis en separar el correo no deseado publicitario (\textit{spam}) de los mensajes fraudulentos diseñados para el engaño (\textit{smishing}).

A lo largo del capítulo se describe cómo se han agrupado trece fuentes distintas, cómo se han resuelto los conflictos de etiquetas entre ellas, y cómo se ha normalizado el texto aplicando técnicas de anonimización profunda y deduplicación. Finalmente, se detalla el enriquecimiento semántico mediante Modelos de Lenguaje de Gran Escala (\glspl{LLM}) e inteligencia de dominios (WHOIS).

\section{Agregación y consolidación de fuentes}
\label{sec:recopilacion}

Los conjuntos de datos (\textit{datasets}) públicos de SMS presentan multitud de formatos, idiomas y formas de etiquetar la información. Para superar este problema y evitar que el modelo se acostumbre a un único entorno, hemos diseñado un flujo de trabajo (\textit{pipeline}) que unifica múltiples repositorios bajo un mismo esquema.

\subsection{Diversidad y procedencia de los datos}
El corpus maestro integra trece repositorios diferentes, seleccionados para abarcar tanto las amenazas históricas como las más recientes:
\begin{itemize}
    \item \textbf{Bases de referencia clásicas}: \textit{UCI SMS Spam Collection} y el \textit{NUS SMS Corpus}, fundamentales para establecer la línea base de los correos legítimos (\textit{ham}) y del \textit{spam} tradicional.
    \item \textbf{Aportaciones recientes (2022-2025)}: Colecciones como \textit{Mishra \& Soni (2022)}, \textit{Hosseinpour (2025)} y el conjunto de \textit{Agarwal IMC (2025)}. Incluirlas garantiza que el modelo aprenda de las campañas de \textit{smishing} actuales.
    \item \textbf{Corpus multiclase y especializados}: \textit{MIMICS-3500} y \textit{Smishing-4C}, que aportan categorías detalladas sobre técnicas de persuasión.
    \item \textbf{Colecciones en español y locales}: Dado el contexto de la investigación, es crítico incluir conjuntos como \textit{Spanish Spam Ham} y extracciones directas de copias de seguridad de dispositivos Android (\textit{ExAIS SMS}).
    \item \textbf{Otras agrupaciones}: Diversos corpus alojados en plataformas de ciencia de datos (\textit{Kaggle Phishing}, \textit{Kaggle Spam/Ham}) y el conjunto \textit{Malicious/Benign SMS/MMS}.
\end{itemize}

Ingerir esta enorme variedad de repositorios no es un proceso trivial. Para resolverlo, el código fuente orquesta un ecosistema de trece adaptadores especializados (\textit{Loaders}). Estos módulos actúan como traductores dinámicos: son capaces de extraer información de jerarquías complejas en formato JSON (como ocurre en el \textit{NUS SMS Corpus}), mapear tabuladores anidados en archivos TSV (\textit{UCI}), e incluso reconstruir lógicamente mensajes fraccionados en múltiples celdas CSV tras ser exportados en crudo desde dispositivos móviles (\textit{ExAIS}). Estos adaptadores aseguran que toda la ingesta masiva converja hacia una matriz de datos estructuralmente idéntica antes de llegar a la etapa de limpieza.

\subsection{Armonización de etiquetas y resolución de conflictos}
\label{subsec:armonizacion}
Dada la variedad de etiquetas originales (términos como \textit{legitimate}, \textit{promotional}, \textit{phishing} o \textit{fraud}), es necesario establecer un mapeo claro hacia nuestras tres clases objetivo: \texttt{ham}, \texttt{spam} y \texttt{smishing}.

Este proceso requiere reglas muy precisas cuando las fuentes presentan ambigüedades. Por ejemplo, el conjunto de \textit{Hosseinpour (2025)} utiliza dos marcas binarias independientes (\texttt{spam\_label} y \texttt{smishing\_label}). En nuestro sistema, hemos establecido una prioridad basada en la severidad de la amenaza: si \texttt{smishing\_label == 1}, la muestra se cataloga directamente como \texttt{smishing}; si no lo es, pero tiene \texttt{spam\_label == 1}, se etiqueta como \texttt{spam}. En fuentes más simples como \textit{Kaggle Phishing}, se hace un mapeo directo hacia \texttt{smishing}.

Para asegurar que todo el proceso se pueda reproducir y auditar, el dataset final siempre conserva la etiqueta original (\texttt{original\_label}) y registra la regla exacta que se usó para transformarla (\texttt{label\_mapping\_rule}).

\section{Limpieza, Anonimización y Deduplicación}
\label{sec:limpieza_anonimizacion}

Esta fase garantiza que los datos sean seguros de procesar, no contengan duplicados y eviten cualquier sesgo circular durante el entrenamiento.

\subsection{Filtro anti-IA: Exclusión de texto generado automáticamente}
Es vital evitar que el corpus de entrenamiento se contamine con mensajes creados por otros LLMs (una práctica que a veces se usa para inflar los datasets artificialmente). Si entrenáramos el modelo con estos datos, correríamos el riesgo de que aprendiera a identificar el estilo de redacción de una inteligencia artificial, en lugar de captar la intención maliciosa de un humano. Esto daría una falsa sensación de precisión (\textit{data leakage} estilístico). Por eso, descartamos explícitamente cualquier muestra que venga marcada de origen como generada por IA (\texttt{is\_ai\_generated}), protegiendo así la naturaleza orgánica del corpus.

\subsection{Limpieza textual y marcas de ofuscación}
\label{subsec:limpieza_ofuscacion}

El preprocesamiento del corpus MetaSMS-HSS se realiza mediante el módulo determinista \texttt{DatasetCleaner} (ubicado en \texttt{utils/cleaner.py}). La justificación metodológica de este componente no radica únicamente en la tradicional limpieza de ruido tipográfico, sino en resolver de manera activa dos problemas críticos: el \textit{data leakage} derivado de la fusión de datasets heterogéneos y la detección explícita de técnicas de evasión empleadas por atacantes reales.

\subsubsection{Prevención de Data Leakage Estilístico por Tokens Híbridos}
Al integrar trece fuentes académicas diferentes, se observa que los autores de cada corpus emplean esquemas de anonimización muy diversos. Por ejemplo, el dataset A sustituye enlaces mediante \texttt{[URL]} o \texttt{<LINK>}, el dataset B enmascara teléfonos usando la cadena ficticia \texttt{0000000000} o el token \texttt{[PHONE]}, mientras que el dataset C mantiene texto crudo. Si un clasificador de aprendizaje automático (\textit{Machine Learning}) se entrena con esta disparidad, se produce un sesgo de origen catastrófico: el clasificador aprende a memorizar que el token estilístico \texttt{<LINK>} está fuertemente correlacionado con la clase \texttt{ham} (porque el dataset legítimo origen lo usaba) y que los enlaces crudos o tokens como \texttt{<URL\_HTTPS>} pertenecen al \texttt{smishing}. 

Para neutralizar este factor de confusión, el motor ejecuta un sub-módulo de estandarización léxica estricta antes de cualquier otra transformación. Esto unifica todas las representaciones placeholders bajo una misma ontología de tokens: \texttt{<URL>}, \texttt{<PHONE>}, \texttt{<EMAIL>}, \texttt{<PERSON>}, \texttt{<CREDIT\_CARD>}, \texttt{<ACCOUNT\_NUMBER>} y \texttt{<DATE>}.

\subsubsection{Mitigación y Codificación de 6 Patrones de Ofuscación}
Para evadir las pasarelas de seguridad (antispam) tradicionales basadas en firmas y palabras clave, las campañas de fraude inyectan distorsiones estructuradas en el texto. \texttt{DatasetCleaner} caza y neutraliza activamente seis variantes de evasión mediante los siguientes mecanismos:
\begin{enumerate}
    \item \textbf{Caracteres Invisibles o de Control (\textit{Zero-Width Spaces})}: Uso de secuencias no imprimibles como espacios de ancho cero (\texttt{\textbackslash u200B-\textbackslash u200F}), modificadores bidireccionales (\texttt{\textbackslash u202A-\textbackslash u202E}), marcas de orden de bytes (\texttt{\textbackslash uFEFF}) y guiones suaves (\texttt{\textbackslash u00AD}). Estos caracteres interrumpen la correspondencia de palabras clave a nivel de bytes (ej. \texttt{B\textbackslash u200Ba\textbackslash u200Bn\textbackslash u200Bc\textbackslash u200Bo}) pero son invisibles en pantalla. El limpiador los elimina y registra la marca \texttt{<OBF\_ZWSP>}.
    \item \textbf{Homoglifos Multilingües (Cyrillic/Greek Confusables)}: Mapeo de caracteres cirílicos o griegos que comparten idéntico glifo visual con el alfabeto latino (por ejemplo, la ``а'' cirílica [\texttt{\textbackslash u0430}] o la ``ο'' griega [\texttt{\textbackslash u03bf}] en lugar de la ``a'' y ``o'' latinas). El sistema utiliza un diccionario de equivalencias estático y traduce cada carácter al bloque latino estándar para su normalización semántica, etiquetando la anomalía con \texttt{<OBF\_HOMOGLYPH>}.
    \item \textbf{Ofuscación por Espaciado (Spacing)}: Inyección artificial de espacios en blanco entre cada una de las letras de un término crítico (como \texttt{C O R R E O S}). Se detecta a través de la expresión regular \texttt{\textbackslash b(?:[a-zA-Z]\textbackslash s)\{3,\}[a-zA-Z]\textbackslash b}, colapsando la palabra afectada y catalogándola como \texttt{<OBF\_SPACING>}.
    \item \textbf{Interrupción por Signos de Puntuación (Punctuation)}: Separación de caracteres mediante puntos, guiones o barras bajas (por ejemplo, \texttt{P.A.Q.U.E.T.E} o \texttt{c\_o\_r\_r\_e\_o\_s}). Se localiza con el patrón \texttt{\textbackslash b(?:[a-zA-Z][.\textbackslash -\_])\{3,\}[a-zA-Z]\textbackslash b}, procediendo a purgar los caracteres de puntuación internos y añadiendo el metadato \texttt{<OBF\_PUNCTUATION>}.
    \item \textbf{Saltos de Línea Excesivos (Newline Evasion)}: Inserción de múltiples saltos de línea consecutivos concebidos para empujar los enlaces fraudulentos fuera de la pantalla del terminal móvil (forzando al usuario a realizar \textit{scroll} y ocultando el hipervínculo a simple vista). Se colapsan todas las repeticiones de espacios, tabulaciones y retornos de carro a un único espacio plano, y si se detecta una secuencia de dos o más retornos de carro, se etiqueta como \texttt{<OBF\_NEWLINE>}.
    \item \textbf{Normalización Tipográfica (Unicode Math Fonts)}: Uso de bloques de fuentes matemáticas en negrita, cursiva o de estilo gótico (como \texttt{𝕻𝖔𝖘𝖙𝖆𝖑}) o caracteres de ancho completo (\texttt{ｅ}). El limpiador aplica la forma de normalización de compatibilidad Unicode NFKC (\textit{Normalization Form Compatibility Decomposition}) para aplanar estos caracteres a ASCII, asignando la marca \texttt{<OBF\_FONT>} en caso de registrarse alteraciones en el texto.
\end{enumerate}

Es fundamental resaltar que el limpiador \textbf{no concatena físicamente las marcas de ofuscación al final de la cadena de texto de entrada}, sino que las almacena de forma aislada en un campo estructurado de metadatos (\texttt{obfuscation\_tags}). Esto evita contaminar sintácticamente el texto plano y permite que los modelos de Deep Learning (como DeBERTa) incorporen estas marcas de manera desacoplada en su cabezal de clasificación o canal de embedding.

\subsubsection{Heurística de Validación de Formato SMS}
A fin de evitar la infiltración de correos electrónicos, boletines extensos o textos descriptivos procedentes de las exportaciones en crudo de los datasets, el limpiador implementa el método heurístico \texttt{is\_valid\_sms}. Este excluye de manera fulminante cualquier mensaje que supere el límite físico extendido de 1600 caracteres o que incorpore firmas típicas de correos electrónicos (tales como \textit{``click here to unsubscribe''}, \textit{``view email in browser''}, o \textit{``all rights reserved''}).

El código completo que orquesta esta normalización textual y extracción de marcas se expone en el Listado~\ref{lst:cleaner_py}.

\begin{lstlisting}[language=Python, caption={Método de limpieza y normalización textual en DatasetCleaner.}, label=lst:cleaner_py]
def clean(self, text: str) -> Dict[str, Any]:
    if not isinstance(text, str):
        text = str(text)
        
    self.stats["total_processed"] += 1
    tags = []
    original = text

    # Normalizar secuencias de escape de espacios en blanco de exportaciones
    if "\\n" in text or "\\t" in text or "\\r" in text:
        text = text.replace("\\r", " ").replace("\\n", " ").replace("\\t", " ")

    # 0. Estandarizar tokens de pre-anonimizacion de datasets academicos
    text = re.sub(r'(?i)<URL>|\[URL\]|http://url\.com|https://url\.com|http://link\.com|https://link\.com|<LINK>|\[LINK\]', '<URL>', text)
    text = re.sub(r'(?i)<PHONE>|\[PHONE\]|<PHONE_NUMBER>|\[PHONE_NUMBER\]|\b0000000000\b', '<PHONE>', text)
    text = re.sub(r'(?i)<EMAIL>|\[EMAIL\]', '<EMAIL>', text)
    text = re.sub(r'(?i)<USER>|\[USER\]|<USERNAME>|\[USERNAME\]|<NAME>|\[NAME\]|<NAMED_ENTITY>|\[NAMED_ENTITY\]', '<PERSON>', text)
    text = re.sub(r'(?i)<CARD>|\[CARD\]|<CREDIT_CARD>|\[CREDIT_CARD\]', '<CREDIT_CARD>', text)
    text = re.sub(r'(?i)<ACCOUNT>|\[ACCOUNT\]|<ACC>|\[ACC\]', '<ACCOUNT_NUMBER>', text)
    text = re.sub(r'(?i)<DATE_TIME>|\[DATE_TIME\]|<DATETIME>|\[DATETIME\]', '<DATE>', text)

    # 1. Detectar y eliminar caracteres invisibles (ZWSP)
    if self.invisible_chars_re.search(text):
        tags.append("<OBF_ZWSP>")
        self.stats["obf_zwsp"] += 1
        text = self.invisible_chars_re.sub('', text)

    # 2. Detectar y sustituir homoglifos
    if self.homoglyph_re.search(text):
        tags.append("<OBF_HOMOGLYPH>")
        self.stats["obf_homoglyph"] += 1
        trans_table = str.maketrans(self.HOMOGLYPH_MAP)
        text = text.translate(trans_table)

    # 3. Detectar ofuscacion por espaciado (C O R R E O S)
    if self.spacing_re.search(text):
        tags.append("<OBF_SPACING>")
        self.stats["obf_spacing"] += 1
        text = self.spacing_re.sub(lambda m: m.group(0).replace(' ', ''), text)
        
    # 4. Detectar ofuscacion por puntuacion (P.A.Q.U.E.T.E)
    if self.punct_obf_re.search(text):
        tags.append("<OBF_PUNCTUATION>")
        self.stats["obf_punct"] += 1
        text = self.punct_obf_re.sub(lambda m: re.sub(r'[.\-_]', '', m.group(0)), text)

    # 5. Detectar saltos de linea excesivos
    newline_matches = self.newlines_re.findall(text)
    if newline_matches:
        if any(len(match) >= 2 for match in newline_matches):
            tags.append("<OBF_NEWLINE>")
            self.stats["obf_newlines"] += 1
        text = re.sub(r'[\r\n\t]+', ' ', text).strip()

    if "\t" in text:
        text = re.sub(r"\t+", " ", text).strip()

    # 6. Normalizacion Unicode (NFKC) para fuentes matematicas
    normalized_text = unicodedata.normalize("NFKC", text)
    if normalized_text != text:
        tags.append("<OBF_FONT>")
        self.stats["obf_font"] += 1
        text = normalized_text

    return {
        "text_raw": original,
        "text_normalized": text,
        "tags": tags
    }
\end{lstlisting}

\subsection{Filtro de privacidad estricta (PII)}
\label{subsec:filtro_privacidad}

La anonimización de Información Personal Identificable (PII) en MetaSMS-HSS está encomendada al módulo \texttt{SMSAnonymizer} (\texttt{utils/anonymizer.py}). El cumplimiento normativo del RGPD exige un borrado seguro de identificadores personales, pero desde el punto de vista del modelado predictivo, la anonimización cumple un rol fundamental de abstracción semántica: obliga al clasificador a enfocarse en la \textit{estructura psicológica} de la ingeniería social en lugar de correlacionar de forma sesgada dominios web específicos, números telefónicos concretos, o identificadores temporales que caducarán en el mundo real en cuestión de días.

\subsubsection{Arquitectura Híbrida de Detección}
Para optimizar el balance entre precisión de detección y velocidad de cómputo, el módulo implementa una arquitectura en cascada de doble motor:
\begin{itemize}
    \item \textbf{Motor Algorítmico por Expresiones Regulares}: Un conjunto exhaustivo de expresiones estructuradas y ordenadas de forma descendente, desde las más específicas a las más genéricas, para evitar solapamientos parciales no deseados.
    \item \textbf{Motor Neuronal de Clasificación de Tokens}: Integración opcional de un pipeline de HuggingFace optimizado para la detección de entidades de privacidad (\texttt{openai/privacy-filter}), ejecutado de forma asíncrona sobre hardware de GPU cuando está habilitado.
\end{itemize}

Antes de ejecutar el escaneo, el texto en crudo experimenta un preprocesamiento crítico que neutraliza las estrategias de evasión de PII: decodifica las direcciones y parámetros utilizando \texttt{urllib.parse.unquote} (para asegurar que las URLs ofuscadas mediante codificación por porcentaje, tipo \texttt{\%3A\%2F\%2F}, sean capturadas) y realiza una normalización unicode NFKC, purgando caracteres no imprimibles del rango ASCII (\texttt{[\textbackslash x00-\textbackslash x08\textbackslash x0B\textbackslash x0C\textbackslash x0E-\textbackslash x1F\textbackslash x7F]}).

\subsubsection{Resolución de Colisiones y Reglas de Solapamiento}
Cuando conviven múltiples motores de detección, surge el problema de solapamiento de límites (un mismo fragmento de texto es marcado parcialmente por Regex y de manera distinta por el clasificador neuronal). Para asegurar una sustitución atómica, el motor implementa las siguientes heurísticas de decisión:
\begin{enumerate}
    \item \textbf{Prioridad de Origen}: Se prioriza de forma absoluta las marcas del motor Regex (Prioridad 1) sobre las detecciones neuronales (Prioridad 0), garantizando que patrones estructurados (como IBANs, tarjetas o claves criptográficas) se enmascaren con tokens ultra-específicos antes de delegar a la inferencia probabilística del LLM/Transformers.
    \item \textbf{Criterio de Longitud de Span}: Ante conflicto entre spans de idéntica prioridad, se retiene el de mayor longitud de cobertura física sobre el string, descartando las detecciones cortas o anidadas.
    \item \textbf{Reemplazo de Derecha a Izquierda}: Al sustituir los fragmentos por marcas fijas (como \texttt{<CRYPTO\_ADDRESS>} en lugar de una dirección de 40 caracteres), el tamaño total de la cadena varía dinámicamente. Para evitar la desalineación de los índices posicionales calculados en la fase de escaneo, las operaciones de modificación del string se procesan estrictamente en sentido inverso (de derecha a izquierda, ordenando los spans de forma descendente por su índice de inicio).
    \item \textbf{Deduplicación de Metamarcas}: Tras el reemplazo, si se producen marcas repetidas de forma consecutiva debido a fragmentaciones del texto (como \texttt{<URL><URL>}), estas se colapsan dinámicamente a un único token empleando expresiones regulares.
\end{enumerate}

El Listado~\ref{lst:anonymizer_py} muestra la lógica de compilación de spans y el proceso de enmascaramiento aséptico en lotes del tokenizador de privacidad.

\begin{lstlisting}[language=Python, caption={Método de procesamiento de mensajes en SMSAnonymizer.}, label=lst:anonymizer_py]
def process_messages(self, texts: List[str]) -> List[Dict[str, Any]]:
    processed_texts = []
    for text in texts:
        if not isinstance(text, str):
            text = str(text)
        text = self._decode_urls(text)
        text = self._normalize_unicode(text)
        text = re.sub(r"[\x00-\x08\x0B\x0C\x0E-\x1F\x7F]", "", text)
        processed_texts.append(text)

    entities_batch = [[] for _ in texts]
    if self.use_privacy_filter and self._privacy_pipe:
        entities_batch = self._process_with_privacy_filter_batch(processed_texts)

    results = []
    for i, text in enumerate(processed_texts):
        self.stats["total"] += 1
        all_spans: List[Tuple[int, int, str, int]] = []
        
        # 1. Spans del motor Regex (Prioridad 1)
        for pattern, replacement, *flags in self.patterns:
            flag = flags[0] if flags else 0
            for match in re.finditer(pattern, text, flag):
                start, end = match.start(1), match.end(1)
                all_spans.append((start, end, replacement, 1))
                
        # 2. Spans del motor de Inteligencia Artificial (Prioridad 0)
        if self.use_privacy_filter and self._privacy_pipe:
            ai_output = entities_batch[i]
            for entity in ai_output:
                label = entity["entity_group"]
                normalized = self._normalize_privacy_label(label)
                token = self.PRIVACY_LABEL_MAP.get(normalized)
                if token:
                    all_spans.append((entity["start"], entity["end"], token, 0))

        # 3. Resolucion de solapamientos: prioridad Regex, luego longitud
        all_spans.sort(key=lambda x: (x[3], (x[1] - x[0]), -x[0]), reverse=True)
        
        final_spans = []
        for start, end, replacement, priority in all_spans:
            overlap = False
            for s, e, _, _ in final_spans:
                if max(start, s) < min(end, e):
                    overlap = True
                    break
            if not overlap:
                final_spans.append((start, end, replacement, priority))
                
        # 4. Sustitucion ordenada de derecha a izquierda
        final_spans.sort(key=lambda x: x[0], reverse=True)
        
        anon_text = text
        for start, end, replacement, priority in final_spans:
            self.stats["replaced"] += 1
            anon_text = anon_text[:start] + replacement + anon_text[end:]

        # Deduplicar marcas consecutivas (ej. <URL> <URL>)
        anon_text = re.sub(r"(<[a-zA-Z0-9_]+>)(?:\s*\1)+", r"\1", anon_text)

        results.append({
            "original": text,
            "anonymized": anon_text,
            "entities": entities_batch[i] if (self.use_privacy_filter and self._privacy_pipe) else []
        })

    return results
\end{lstlisting}

\subsection{Deduplicación cruzada con SHA-256}
\label{subsec:deduplicacion}

Al agregar múltiples fuentes públicas de ciberseguridad, es sumamente frecuente que un mismo mensaje fraudulento o legítimo se encuentre repetido en varios repositorios. Si no se filtra con rigor, un mismo mensaje podría distribuirse simultáneamente en los conjuntos de entrenamiento y de validación/testeo, provocando un sobreajuste artificial (\textit{overfitting}) y distorsionando la validez estadística de las métricas de rendimiento.

Para conjurar este peligro, el pipeline calcula una firma criptográfica única mediante el algoritmo SHA-256 sobre una versión normalizada del texto. El reto metodológico consiste en evitar que diferencias insignificantes impidan asociar dos mensajes idénticos. Por ejemplo, un mismo SMS puede haber sido registrado por una fuente con marcas de ofuscación (ej. ZWSP) y por otra en texto claro. Para corregir esto, antes de computar el hash, el sistema elimina mediante expresiones regulares cualquier etiqueta de ofuscación inyectada previamente por el limpiador textual. Posteriormente, se aplica la normalización Unicode NFKC, se convierte la cadena a minúsculas y se colapsa todo espacio en blanco. El Listado~\ref{lst:dedupe} muestra la función implementada para este cálculo.

\begin{lstlisting}[language=Python, caption={Cálculo de la firma criptográfica para deduplicación cruzada.}, label=lst:dedupe]
def build_dedupe_key(text: str) -> str:
    """Hash a normalized text string for cross-source deduplication.

    The input should be the cleaned+anonymized text without OBF annotation
    tags so that the same message appearing in two datasets
    produces the same hash.
    """
    if not text:
        return ""
    # Strip OBF tags appended by DatasetCleaner before hashing.
    clean = re.sub(r"\s*<OBF_[A-Z_]+>", "", text)
    normalized = unicodedata.normalize("NFKC", clean)
    normalized = " ".join(normalized.lower().split())
    return hashlib.sha256(normalized.encode("utf-8", errors="ignore")).hexdigest()
\end{lstlisting}

\section{Esquema Unificado y Metadatos}
\label{sec:esquema_trazabilidad}

Al terminar la primera fase de construcción, toda la información se estructura bajo el esquema base, dividido en tres bloques estandarizados (Tablas~\ref{tab:esquema_core}, \ref{tab:esquema_trazabilidad} y \ref{tab:esquema_linguistico}).

\begin{table}[htbp]
    \centering
    \caption{Esquema MetaSMS-HSS: Datos básicos y clases.}
    \label{tab:esquema_core}
    \small
    \begin{tabular}{llp{7cm}}
        \toprule
        \textbf{Campo} & \textbf{Tipo} & \textbf{Descripción} \\
        \midrule
        \texttt{message\_id} & String (UUID) & Identificador hexadecimal único interno. \\
        \texttt{text\_raw} & String & Contenido textual original en bruto sin procesar (con posibles ofuscaciones y PII). \\
        \texttt{text} & String & Contenido textual normalizado tras la mitigación de los 6 patrones de ofuscación (texto procesado). \\
        \texttt{text\_anonymized} & String & Contenido textual tras aplicar el motor de anonimización híbrido de privacidad (PII). \\
        \texttt{canonical\_label} & Categorical & Variable objetivo a predecir: \texttt{ham}, \texttt{spam}, \texttt{smishing}. \\
        \bottomrule
    \end{tabular}
\end{table}

\begin{table}[htbp]
    \centering
    \caption{Esquema MetaSMS-HSS: Procedencia y auditoría.}
    \label{tab:esquema_trazabilidad}
    \small
    \begin{tabular}{llp{7cm}}
        \toprule
        \textbf{Campo} & \textbf{Tipo} & \textbf{Descripción} \\
        \midrule
        \texttt{reference} & String & Nombre de la fuente original. \\
        \texttt{source\_id} & String & Identificador principal en el repositorio de origen. \\
        \texttt{original\_label} & String & Etiqueta original antes de la armonización. \\
        \texttt{label\_mapping\_rule} & String & Regla de transformación que se aplicó a la etiqueta. \\
        \texttt{timestamp\_original} & Datetime / NULL & Marca de tiempo original de recepción (si existiera). \\
        \texttt{license} & String & Licencia de uso original. \\
        \bottomrule
    \end{tabular}
\end{table}

La asignación del idioma (\texttt{language}) plantea un reto técnico de gran calado empírico debido a la brevedad y el ruido tipográfico del formato SMS. Para solventarlo de raíz de manera eficiente, el \textit{pipeline} implementa un motor basado en la red neuronal convolucional \textbf{FastText de Meta} (empleando el modelo comprimido \texttt{lid.176.bin}). El sistema procesa el texto en ráfaga infiriendo el idioma y su nivel de confianza probabilística. Como técnica de contingencia ante la naturaleza informal del SMS, si la confianza inicial es inferior al $50\%$ y el mensaje contiene letras mayúsculas, el motor realiza un segundo intento forzando el texto a minúsculas, lo que mitiga los falsos negativos provocados por el énfasis tipográfico del atacante. Finalmente, los códigos de idioma resultantes se estandarizan a su representación ISO 639-1 de dos letras utilizando la librería \texttt{pycountry}. Este enfoque asegura una catalogación lingüística robusta sin la sobrecarga computacional de motores probabilísticos tradicionales.

\begin{table}[htbp]
    \centering
    \caption{Esquema MetaSMS-HSS: Procesamiento lingüístico.}
    \label{tab:esquema_linguistico}
    \small
    \begin{tabular}{llp{7cm}}
        \toprule
        \textbf{Campo} & \textbf{Tipo} & \textbf{Descripción} \\
        \midrule
        \texttt{language} & String & Código ISO 639-1 inferido por el motor jerárquico. \\
        \texttt{language\_confidence} & Float & Nivel de confianza probabilística del modelo ganador. \\
        \texttt{anonymization\_status} & Categorical & Estado de alteración estructural por privacidad. \\
        \bottomrule
    \end{tabular}
\end{table}

\section{Enriquecimiento de Inteligencia (IA y WHOIS)}
\label{sec:enriquecimiento}

La fase de enriquecimiento inyecta contexto externo a los mensajes. Estos metadatos son extremadamente valiosos para aislar comportamientos durante la fase de investigación, pero no se utilizan como variables predictoras de entrada en el modelo final para evitar corrupciones de rendimiento.

\subsection{Anotación contextual con LLMs: Arquitectura Local vLLM de Alto Rendimiento}
\label{subsec:anotacion_llm}

Para analizar de forma cualitativa y semántica la intención de los mensajes, el pipeline incorpora un módulo de enriquecimiento semántico mediante Modelos de Lenguaje de Gran Escala (\glspl{LLM}). A diferencia de los enfoques tradicionales que dependen de APIs comerciales en la nube (con los consiguientes riesgos de privacidad, latencia y costes de facturación), nuestra arquitectura ejecuta de forma local el modelo cuantizado de 35 billones de parámetros (Mixture of Experts con 3B activos) **Qwen/Qwen3.5-35B-A3B-GPTQ-Int4** mediante la librería de optimización de inferencia **vLLM** sobre el hardware GPU NVIDIA A100 (40GB) del clúster de supercomputación Legio.

\subsubsection{Mitigación de Cuellos de Botella y Optimización HPC}
La inferencia de modelos a gran escala sobre corpus masivos presenta importantes retos de ingeniería de datos. Para evitar la lentitud que comúnmente penaliza estas ejecuciones, se implementaron cuatro optimizaciones críticas para neutralizar los cuellos de botella de hardware y software:
\begin{enumerate}
    \item \textbf{Eliminación del Cuello de Botella I/O de Red (CEPH vs. Local SSD)}: El manual del clúster Legio advierte explícitamente que la red de almacenamiento en red CEPH (donde se ubica la carpeta \texttt{/home}) penaliza severamente el acceso a múltiples archivos pequeños. Para solventar este retardo, el flujo de trabajo copia el dataset y los archivos intermedios directamente al disco local de estado sólido SSD del nodo computacional (\texttt{/local\_scratch}), reduciendo la latencia de acceso a cero y garantizando que la GPU no pase tiempo inactiva esperando datos.
    \item \textbf{Inferencia en Lotes Masivos (Batching)}: En lugar de procesar los mensajes de forma secuencial individualizada (tamaño de lote igual a 1), se configuró una ejecución con un tamaño de lote (\texttt{batch\_size}) de 128 mensajes. Esto satura el bus de la GPU, permitiendo que la arquitectura paralela procese las muestras simultáneamente y aumente el rendimiento hasta 50 veces.
    \item \textbf{Inferencia Optimizada mediante vLLM}: El uso del motor estándar de Hugging Face (\texttt{model.generate()}) presenta una gestión ineficiente de la memoria RAM de vídeo (VRAM). Para blindar la ejecución frente a errores de falta de memoria (\textit{Out Of Memory} - OOM), se integró la librería **vLLM**, la cual aprovecha las técnicas de **PagedAttention** y **Continuous Batching** para optimizar el direccionamiento de memoria de claves y valores (KV cache), maximizando el throughput del clúster.
    \item \textbf{Decodificación Guiada (Guided JSON Decoding)}: Para evitar la latencia lineal derivada de la generación de respuestas extensas o el razonamiento no estructurado del modelo, se inyecta un esquema estricto de decodificación guiada (\texttt{guided\_json}). Esto fuerza al modelo a restringir sintácticamente su generación a una estructura JSON y detener la inferencia inmediatamente tras cerrar el objeto, limitando el volumen de salida a unos 40--50 tokens por mensaje.
\end{enumerate}

Para cada mensaje, el LLM recibe una directiva de sistema (\textit{system prompt}) estricta que le fuerza a devolver un objeto JSON estructurado con las siguientes claves:
\begin{itemize}
    \item \texttt{theme}: Temática o intención principal del SMS, restringida a una taxonomía cerrada mutuamente excluyente y colectivamente exhaustiva (\textit{MECE}) de 14 categorías: \texttt{personal}, \texttt{family\_emergency}, \texttt{banking}, \texttt{delivery}, \texttt{account\_security}, \texttt{tech\_support}, \texttt{promotion}, \texttt{survey}, \texttt{government}, \texttt{job\_offer}, \texttt{dating\_adult}, \texttt{subscription}, \texttt{service\_alert} y \texttt{unknown}.
    \item \texttt{urgency\_level}: Nivel de presión psicológica discreto que ejerce el mensaje sobre el receptor, categorizado en: \texttt{none}, \texttt{low}, \texttt{medium} o \texttt{high}.
    \item \texttt{language}: Código ISO 639-1 del idioma principal detectado semánticamente en el texto.
\end{itemize}

\subsubsection{Heurística de Fallback Lingüístico}
La integración del LLM aporta además un mecanismo de contingencia para mejorar la detección lingüística del corpus. Aunque el motor de FastText es sumamente veloz, su precisión decae en mensajes con una alta densidad de jerga regional (por ejemplo, el \textit{Singlish} de Singapur o abreviaciones extremas en SMS). Para corregir esto, el script de enriquecimiento (\texttt{enrich\_metadataset.py}) evalúa la confianza reportada por FastText: si esta es inferior al $60\%$ o si el idioma fue catalogado como indeterminado (\texttt{unknown}), el pipeline sobrescribe el idioma utilizando la clasificación semántica realizada por el modelo Qwen, asignando una confianza especial con el token literal \texttt{"LLM"}.

\subsection{Inteligencia WHOIS y heurística de puntuación}
Dado que el enlace fraudulento constituye el núcleo de la ingeniería social en un \textit{smishing}, el módulo \texttt{WhoisEnricher} intercepta las URLs del SMS antes de que la capa de privacidad las destruya. A continuación, contacta a la red registral global de Internet para interrogar el dominio subyacente.

Suele ocurrir que un mismo ataque inyecta \textbf{múltiples enlaces simultáneos} (por ejemplo, camuflando una URL falsa entre enlaces legítimos de descarga de la aplicación bancaria real para cimentar su credibilidad). Para automatizar la resolución del dominio a investigar, se ha programado una heurística matemática de puntuación. Como se observa en el Listado~\ref{lst:whois}, el algoritmo recompensa a los dominios que blindan la identidad del propietario mediante *proxys* comerciales y los penaliza duramente según los días que hayan transcurrido desde su creación en Internet. 

\begin{lstlisting}[language=Python, caption={Algoritmo heurístico para seleccionar el dominio sospechoso principal en un SMS.}, label=lst:whois]
best_w_data = None
best_score = -9999

for url in set(urls):
    w_data = whois_enricher.get_whois_data(url)
    if w_data:
        # Puntuacion base: +10 si utiliza escudo de privacidad comercial
        score = 10 if w_data["hidden"] else 0
        
        # Penalizacion temporal: se restan los dias exactos de antiguedad
        if w_data["age_days"] >= 0:
            score -= w_data["age_days"]
        else:
            score -= 365 # Penalizacion estandar si el servidor esconde la edad
            
        if score > best_score:
            best_score = score
            best_w_data = w_data
\end{lstlisting}

Este dominio ganador se extrae y parametriza directamente como un apéndice más del mensaje, posibilitando posteriormente su análisis estadístico profundo (Tabla~\ref{tab:esquema_llm}).

\begin{table}[htbp]
    \centering
    \caption{Esquema MetaSMS-HSS: Enriquecimiento WHOIS y métricas LLM.}
    \label{tab:esquema_llm}
    \small
    \begin{tabular}{llp{7cm}}
        \toprule
        \textbf{Campo} & \textbf{Tipo} & \textbf{Descripción} \\
        \midrule
        \texttt{llm\_model} & String & Modelo exacto utilizado para la inferencia semántica. \\
        \texttt{theme} & String & Temática principal detectada (ej. cuentas, logística). \\
        \texttt{urgency\_level} & String & Valoración discreta de la presión psicológica ejercida. \\
        \texttt{whois\_domain} & String & Dominio ganador seleccionado por la heurística de riesgo. \\
        \texttt{whois\_age\_days} & Integer / NULL & Días transcurridos desde la creación del registro en Internet. \\
        \texttt{whois\_hidden} & Boolean & Valor 1 si el dominio utiliza protección registral anónima. \\
        \bottomrule
    \end{tabular}
\end{table}

\section{Metodología de Subconjuntos y Muestreo Estratificado}
\label{sec:subconjuntos}
El fichero maestro resultante (con cientos de miles de muestras enriquecidas) no se emplea como un monobloque rígido durante las fases de validación algorítmica. Nuestro sistema de compilación incluye herramientas que permiten aislar particiones deterministas de forma trazable. 

Más allá de simples filtrados por idioma o la supresión de contenido autogenerado por LLMs, el módulo de abstracción (\texttt{create\_subset.py}) orquesta potentes algoritmos de \textbf{muestreo estratificado}. Este requerimiento técnico capacita la extracción instantánea de corpus de experimentación sustancialmente más pequeños manteniendo, invariablemente, las ratios demográficas preexistentes de las tres clases predictoras. El sostenimiento numérico del balanceo original es una precondición innegociable en el campo del \textit{Machine Learning} moderno para sortear sesgos muestrales durante las métricas de rendimiento y la validación cruzada (\textit{Cross-Validation}).

\section{Ejecución Paralelizada en Supercomputación}
\label{sec:ejecucion_hpc}
Todo el *pipeline* exhaustivo descrito en los apartados previos (\textit{hashes} criptográficos iterativos sobre cientos de megabytes, serializaciones JSON con transformadores locales para anonimización y peticiones en ráfaga sostenida a las nubes corporativas de Google o Groq) desborda ampliamente la capacidad térmica y computacional de los equipos de escritorio convencionales.

Por esta razón, la arquitectura del TFM ha sido refactorizada para su procesamiento desatendido en el \textbf{clúster de supercomputación Legio}. La infraestructura está administrada lógicamente por el gestor de recursos SLURM, solicitando despliegues sobre nodos que presenten un hiperparalelismo masivo de CPU y aceleración gráfica subyacente. Para blindar un entorno inmutable que sea reproducible en investigaciones universitarias venideras, todas las librerías transversales del proyecto operan encapsuladas asépticamente mediante contenedores HPC basados en \texttt{cuda-12.4-uv}.

Asimismo, y dado que estas cargas de trabajo asíncronas pueden sostenerse durante un número imprevisible de horas, el sistema subyacente fue diseñado aplicando metodologías estrictas de tolerancia a fallos (\textit{fault-tolerant}). Un motor interno estampa \textbf{puntos de control (\textit{checkpoints})} transaccionales cada 5000 iteraciones, forzando la escritura atómica del estado del procesador en el disco persistente del clúster. Si el nodo se reiniciase por fuerza mayor, superase los tiempos de cómputo autorizados, o las APIs colapsasen, la arquitectura HPC permite que el nodo reanude la compilación en el registro milimétrico donde pereció, evadiendo la necesidad prohibitiva de recalcular algoritmos masivos o consumir cuotas arancelarias redundantes.